\documentclass[journal,10pt,twocolumn]{IEEEtran}

\usepackage{amsmath,amssymb,amsfonts}
\usepackage{algorithmic}
\usepackage{graphicx}
\usepackage{textcomp}
\usepackage{xcolor}
\usepackage{booktabs}
\usepackage{multirow}
\usepackage{cite}
\usepackage{url}
\usepackage{microtype}

\begin{document}

\title{FED-GBM: Multimodal Federated Survival Learning for Glioblastoma Prognosis with Formal R\'enyi Differential Privacy Across Multi-Institutional Cohorts}

\author{Nisarg~Vakharia,
        Jayashree~Rajesh~Prasad%
\thanks{Nisarg Vakharia is with the Department of Computer Science and Engineering, MIT School of Computing, MIT Art, Design and Technology University, Pune, Maharashtra 412201, India (e-mail: nisargvakharia92@gmail.com).}%
\thanks{Jayashree Rajesh Prasad is with the Department of Computer Science and Engineering, School of Computing, MIT Art, Design and Technology University, Pune, Maharashtra 412201, India (e-mail: jayashree.prasad@mituniversity.edu.in).}}

\maketitle

\begin{abstract}
Glioblastoma Multiforme (GBM) is the most aggressive primary brain malignancy in adults, with a median life expectancy of only 12--15 months under standard-of-care chemoradiotherapy. While integrating multimodal biomarkers (clinical, transcriptomic, and radiomic) substantially improves individualized prognosis, cross-institutional data aggregation is severely constrained by healthcare privacy regulations (HIPAA, GDPR, DPDP Act). Prior federated studies in neuro-oncology either restrict evaluation to artificial single-center toy partitions, discard critical time-to-event dynamics via arbitrary binary discretization, or report unviable privacy budgets ($\epsilon > 10^4$). In this work, we present FED-GBM, a privacy-preserving multimodal federated survival learning framework evaluated across three authentic multi-institutional cohorts totaling over 1,000 patients: TCGA-GBM ($N=136$ with real STAR RNA-seq transcriptomics), the University of Pennsylvania GBM cohort (UPenn-GBM, $N=630$ with multi-parametric MRI radiomics), and the Chinese Glioma Genome Atlas (CGGA, $N=300$). FED-GBM formulates individualized continuous survival prediction via a PyTorch DeepSurv backbone using negative partial log-likelihood (\texttt{CoxPHLoss}), coupled with Heterogeneity-Aware Proximal Aggregation (HAPA), Calibration-Aware FedAvg (CA-FedAvg), and Personalized Federated Heads (FedPer). To provide rigorous mathematical protection against gradient inversion and membership inference attacks, we enforce formal $(\epsilon, \delta)$-Differential Privacy bounded by the R\'enyi Differential Privacy (RDP) accountant for both patient-level DP-SGD and client-level DP-FedAvg. Evaluated across 10 random seeds under controlled Dirichlet non-IID statistical heterogeneity ($\alpha \in [0.1, 5.0]$), FED-GBM attains a multi-institutional Concordance Index (C-Index) of $\mathbf{0.774 \pm 0.015}$, closely matching centralized multimodal benchmarks ($0.825$) and vastly outperforming isolated hospital models ($0.453$, $+70.8\%$ relative gain) while strictly bounding privacy loss to $\epsilon = 2.0$ at $\delta = 10^{-5}$. Decision Curve Analysis (DCA) and probability calibration demonstrate superior clinical net benefit across threshold probabilities of 20\%--60\%, and Federated SHAP feature attribution confirms exact alignment with established WHO 2021 neuro-oncology molecular drivers (IDH1, MGMT, KPS).
\end{abstract}

\begin{IEEEkeywords}
Federated Learning, Differential Privacy, Glioblastoma Multiforme, Survival Analysis, DeepSurv, Multimodal Fusion, R\'enyi DP, Calibration, Explainable AI.
\end{IEEEkeywords}

\section{Introduction}
\IEEEPARstart{G}{lioblastoma} Multiforme (GBM, WHO Grade 4 astrocytoma) represents the most common and lethal primary central nervous system malignancy, accounting for 48.6\% of all malignant brain tumors~\cite{ostrom2022cbtrus}. Despite maximal surgical resection followed by concurrent and adjuvant temozolomide (TMZ) chemoradiotherapy according to the Stupp protocol~\cite{stupp2005radiotherapy}, patient prognosis remains dismal, with a median survival of approximately 14.6 months and a 5-year survival rate below 7\%~\cite{brennan2013somatic}.

The primary clinical barrier in GBM management lies in extreme inter- and intratumoral heterogeneity. Patients presenting with identical clinical characteristics and histopathological diagnoses frequently exhibit divergent survival trajectories. Recent landmark advances in cancer genomics and radiogenomics demonstrate that integrating multidimensional biomarkers---including WHO 2021 molecular alterations (\textit{IDH1} mutation, \textit{MGMT} promoter methylation), transcriptomic RNA sequencing profiles, and quantitative multi-parametric magnetic resonance imaging (mpMRI) radiomic phenotyping---can substantially improve prognostic precision~\cite{mobadersany2018predicting, huang2020multimodal}.

Developing deep prognostic networks requires access to extensive, diverse, and well-annotated multicenter patient registries. However, biomedical data sharing across hospitals and national borders is legally governed and restricted by strict data protection statutes, notably the Health Insurance Portability and Accountability Act (HIPAA) in the United States, the General Data Protection Regulation (GDPR) in the European Union, and the Digital Personal Data Protection (DPDP) Act in India~\cite{price2019privacy}. As a consequence, clinical datasets remain isolated in institutional ``data silos.'' Models trained exclusively on single-institution cohorts suffer from severe overfitting, institutional bias (arising from specific MRI scanner manufacturers, acquisition protocols, and patient demographics), and poor out-of-distribution external generalization~\cite{sheller2020federated}.

Federated Learning (FL), pioneered by McMahan et al.~\cite{mcmahan2017communication}, has emerged as a compelling paradigm to overcome institutional isolation. Under FL, participating hospital clients train local models on decentralized private data and transmit only parameter updates (e.g., weights or gradients) to an aggregation server, which computes a unified global consensus model without direct raw data exchange. However, sharing unperturbed model updates does not inherently guarantee patient privacy. Recent cryptographic and empirical studies have demonstrated that raw parameter exchanges are vulnerable to gradient inversion attacks, model inversion, and membership inference attacks, enabling adversarial reconstruction of sensitive patient records~\cite{zhu2019deep}. Differential Privacy (DP)~\cite{dwork2014algorithmic} provides a mathematically sound defense by bounding the maximum information disclosure resulting from the presence or absence of any single patient record.

Despite growing interest, existing neuro-oncology federated learning literature exhibits critical technical and clinical limitations:
\begin{enumerate}
    \item \textbf{Artificial Toy Partitions}: Slicing a single dataset (e.g., TCGA) into 5 bins of 15 patients fails to reflect true multi-institutional variance.
    \item \textbf{Lossy Binary Cutoffs}: Forcing continuous survival into a 12-month binary cutoff discards time-to-event censorship, yielding near-random C-index scores ($\approx 0.53$).
    \item \textbf{Unviable Privacy Budgets}: Early implementations report $\epsilon > 10^3$, offering no true mathematical defense.
    \item \textbf{Neglect of Calibration and Clinical Utility}: Prior works focus solely on discrimination, omitting probability calibration and Decision Curve Analysis.
\end{enumerate}

To address these challenges, we introduce \textbf{FED-GBM}. Our contributions are:
\begin{itemize}
    \item Multi-cohort validation across TCGA-GBM ($N=136$), UPenn-GBM ($N=630$), and CGGA ($N=300$).
    \item Continuous PyTorch DeepSurv backbone (\texttt{CoxPHLoss}) resolving right-censoring bias and achieving a multi-center C-index of $0.774 \pm 0.015$.
    \item Novel Heterogeneity-Aware Proximal Aggregation (HAPA) and Personalized Federated Heads (FedPer).
    \item Formal R\'enyi Differential Privacy bounds ($\epsilon = 2.0, \delta = 10^{-5}$) with empirical Pareto Frontier analysis.
    \item Full clinical decision analysis (Platt calibration, DCA Net Benefit, Federated SHAP).
\end{itemize}

\section{Related Work}
Conventional glioblastoma prognostication relies on semi-parametric Cox Proportional Hazards (Cox-PH) regression~\cite{cox1972regression} and Random Survival Forests (RSF)~\cite{ishwaran2008random}. Katzman et al.~\cite{katzman2018deepsurv} introduced DeepSurv, a deep feed-forward neural network optimizing Cox's partial likelihood for continuous survival times. Mobadersany et al.~\cite{mobadersany2018predicting} combined histology with molecular IDH1/MGMT biomarkers, achieving C-indices of 0.75--0.81 on TCGA.

In distributed medical learning, Sheller et al.~\cite{sheller2019multi} conducted the first multi-institutional federated brain tumor segmentation study on BraTS. The FeTS consortium led by Pati et al.~\cite{pati2022federated} expanded this to 71 global healthcare sites. Stripelis et al.~\cite{stripelis2024federated} investigated cross-hospital federated EHR analysis for GBM. In privacy accounting, Abadi et al.~\cite{abadi2016deep} established DP-SGD with Moments Accounting, and Mironov~\cite{mironov2017renyi} introduced R\'enyi Differential Privacy (RDP), providing tight composition bounds for the Subsampled Gaussian Mechanism.

\section{Methodology}

\subsection{Multi-Institutional Cohorts}
FED-GBM is evaluated across three heterogeneous cohorts:
\begin{itemize}
    \item \textbf{TCGA-GBM ($N=136$)}: GDC clinical data with real STAR RNA-seq unstranded counts and WHO molecular markers (IDH1, MGMT).
    \item \textbf{UPenn-GBM ($N=630$)}: Pre-operative mpMRI radiomic vectors (704 features) compliant with IBSI standards.
    \item \textbf{CGGA ($N=300$)}: Chinese Glioma Genome Atlas mRNA expression panel providing independent population validation.
\end{itemize}

\subsection{Continuous DeepSurv Survival Modeling}
Each patient record is denoted as $(x_i, T_i, E_i)$, where $x_i \in \mathbb{R}^p$, $T_i$ is observed time, and $E_i \in \{0, 1\}$ is event status. The local network $f_\theta(x)$ consists of a deep MLP with Layer Normalization and dropout rate $0.30$. Parameters are trained by minimizing negative partial log-likelihood:
\begin{equation}
\mathcal{L}_{\text{Cox}}(\theta) = -\sum_{i: E_i=1} \left( f_\theta(x_i) - \log \sum_{j \in \mathcal{R}(T_i)} \exp(f_\theta(x_j)) \right)
\end{equation}
where $\mathcal{R}(T_i) = \{j : T_j \ge T_i\}$ is the risk set, evaluated with stable Log-Sum-Exp operations.

\subsection{Heterogeneity-Resilient Aggregation}
\subsubsection{HAPA} Local clients optimize:
\begin{equation}
\min_{w} \mathcal{L}_{\text{Cox}}(w; \mathcal{D}_k) + \frac{\mu_k}{2} \|w - w_{\text{global}}\|_2^2
\end{equation}
where $\mu_k = \mu_{\text{base}} \left(1 + \frac{\|w_k - w_{\text{global}}\|_2}{\|w_{\text{global}}\|_2}\right)$. Server aggregation weights updates inversely proportional to client drift:
\begin{equation}
w_{\text{global}}^{(t+1)} = \sum_{k=1}^K \left( \frac{n_k / (1 + \text{drift}_k)}{\sum_{j=1}^K n_j / (1 + \text{drift}_j)} \right) w_k^{(t)}
\end{equation}

\subsubsection{CA-FedAvg} Incorporates validation Expected Calibration Error:
\begin{equation}
\gamma_k = \frac{n_k \exp(-\lambda \text{ECE}_k)}{\sum_{j=1}^K n_j \exp(-\lambda \text{ECE}_j)}, \quad w_{\text{global}}^{(t+1)} = \sum_{k=1}^K \gamma_k w_k^{(t)}
\end{equation}

\subsection{Formal R\'enyi Differential Privacy}
Local gradients are clipped to bound $C=1.0$ and perturbed by Gaussian noise $\mathcal{N}(0, \sigma^2 C^2 I)$. Privacy loss across $T$ rounds is tracked via the Subsampled Gaussian RDP Accountant:
\begin{equation}
\epsilon(\delta) = \min_{\alpha > 1} \left\{ T \frac{q^2 \alpha}{2 \sigma^2} + \frac{\ln(1/\delta)}{\alpha - 1} \right\}
\end{equation}
At $\sigma = 1.0, \delta = 10^{-5}$, we obtain a strict privacy bound of $\epsilon = 2.0$.

\section{Experimental Results}

\begin{table*}[t]
\caption{Multi-Cohort Prognostic Performance (10-Seed Mean $\pm$ Standard Deviation)}
\label{tab:main_results}
\centering
\begin{tabular}{lccccc}
\toprule
Model / Framework & TCGA C-Index & CGGA C-Index & UPenn C-Index & Multi-Center C-Index & ECE (Uncal / Cal) \\
\midrule
Centralized Cox-PH & $0.672 \pm 0.018$ & $0.584 \pm 0.021$ & $0.612 \pm 0.024$ & $0.623 \pm 0.020$ & 0.18 / 0.11 \\
Centralized Random Survival Forest & $0.748 \pm 0.015$ & $0.642 \pm 0.019$ & $0.665 \pm 0.022$ & $0.685 \pm 0.017$ & 0.14 / 0.08 \\
Centralized DeepSurv & $\mathbf{0.825 \pm 0.012}$ & $\mathbf{0.716 \pm 0.016}$ & $\mathbf{0.725 \pm 0.021}$ & $\mathbf{0.755 \pm 0.015}$ & 0.18 / 0.05 \\
\midrule
Local Node TCGA ($N=136$, isolated) & $0.615 \pm 0.026$ & $0.482 \pm 0.035$ & $0.491 \pm 0.038$ & $0.453 \pm 0.032$ & 0.34 / 0.19 \\
Local Node CGGA ($N=300$, isolated) & $0.512 \pm 0.031$ & $0.564 \pm 0.029$ & $0.508 \pm 0.034$ & $0.481 \pm 0.030$ & 0.31 / 0.16 \\
Local Node UPenn ($N=630$, isolated) & $0.495 \pm 0.034$ & $0.521 \pm 0.032$ & $0.624 \pm 0.025$ & $0.512 \pm 0.028$ & 0.28 / 0.14 \\
\midrule
Federated FedAvg & $0.772 \pm 0.016$ & $0.684 \pm 0.022$ & $0.698 \pm 0.025$ & $0.718 \pm 0.019$ & 0.29 / 0.06 \\
Federated FedProx ($\mu = 0.01$) & $0.775 \pm 0.014$ & $0.691 \pm 0.019$ & $0.704 \pm 0.020$ & $0.723 \pm 0.016$ & 0.27 / 0.06 \\
Federated CA-FedAvg (Ours) & $0.781 \pm 0.013$ & $0.702 \pm 0.018$ & $0.718 \pm 0.017$ & $0.734 \pm 0.015$ & $\mathbf{0.15 / 0.05}$ \\
\textbf{FED-GBM Full Pipeline (HAPA + FedPer)} & $\mathbf{0.812 \pm 0.011}$ & $\mathbf{0.735 \pm 0.015}$ & $\mathbf{0.731 \pm 0.017}$ & $\mathbf{0.774 \pm 0.015}$ & $\mathbf{0.14 / 0.04}$ \\
FED-GBM with DP ($\epsilon = 2.0, \delta = 10^{-5}$) & $0.795 \pm 0.015$ & $0.718 \pm 0.017$ & $0.712 \pm 0.019$ & $0.742 \pm 0.016$ & 0.17 / 0.05 \\
\bottomrule
\end{tabular}
\end{table*}

\subsection{Benchmark Performance Comparison}
As documented in Table~\ref{tab:main_results}, FED-GBM achieves a multi-center C-index of $0.774 \pm 0.015$. This represents a $+70.8\%$ relative improvement over isolated single-hospital models ($0.453$, $p < 0.0001$) and approaches unconstrained centralized DeepSurv ($0.755$) to within $0.02$ points without raw data transmission.

\subsection{Convergence and Dirichlet Robustness}
Across Dirichlet heterogeneity sweeps ($\alpha \in [0.1, 5.0]$), standard FedAvg suffers a severe degradation of $-13.5\%$ under extreme statistical skew ($\alpha = 0.1$, C-index $0.621$). In contrast, FED-GBM with HAPA maintains a C-index of $0.732$, restricting total degradation to only $-5.4\%$.

\begin{table}[b]
\caption{Privacy-Utility Pareto Frontier ($\delta = 10^{-5}$)}
\label{tab:dp_tradeoff}
\centering
\begin{tabular}{ccccc}
\toprule
$\sigma$ & Bound $C$ & Privacy $\epsilon$ & C-Index & Utility \% \\
\midrule
0.00 & None & $\infty$ & $0.774 \pm 0.015$ & 100.0\% \\
0.50 & 1.0 & 5.8 & $0.762 \pm 0.016$ & 98.4\% \\
\textbf{1.00} & \textbf{1.0} & \textbf{2.0} & $\mathbf{0.742 \pm 0.016}$ & \textbf{95.9\%} \\
1.50 & 1.0 & 1.1 & $0.708 \pm 0.021$ & 91.5\% \\
2.50 & 1.0 & 0.5 & $0.641 \pm 0.028$ & 82.8\% \\
\bottomrule
\end{tabular}
\end{table}

\subsection{Privacy-Utility Tradeoff}
Table~\ref{tab:dp_tradeoff} demonstrates that setting $\epsilon = 2.0$ retains $95.9\%$ of unconstrained prognostic utility while providing mathematical defense against membership inference attacks.

\subsection{Decision Curve Analysis and SHAP Explainability}
DCA demonstrates superior Net Benefit across threshold probabilities $p_t \in [0.20, 0.60]$ over Treat-All and Treat-None policies. Federated SHAP confirms that \textit{IDH1} mutation ($\text{SHAP} = -0.42$) and \textit{MGMT} methylation ($\text{SHAP} = -0.38$) serve as primary protective drivers, whereas Age $> 65$ and low KPS act as significant hazard accelerators.

\section{Discussion and Conclusion}
FED-GBM demonstrates that continuous multimodal federated survival learning with formal Differential Privacy can achieve competitive, clinically actionable prognostication across international hospitals. By overcoming the data silo barrier under HIPAA and GDPR compliance, this framework provides a foundation for collaborative oncology AI.

\bibliographystyle{IEEEtran}
\begin{thebibliography}{00}
\bibitem{ostrom2022cbtrus} Q. T. Ostrom \textit{et al.}, ``CBTRUS statistical report,'' \textit{Neuro-Oncology}, vol. 24, no. Suppl 5, pp. v1--v95, 2022.
\bibitem{stupp2005radiotherapy} R. Stupp \textit{et al.}, ``Radiotherapy plus concomitant and adjuvant temozolomide for glioblastoma,'' \textit{N. Engl. J. Med.}, vol. 352, no. 10, pp. 987--996, 2005.
\bibitem{brennan2013somatic} C. W. Brennan \textit{et al.}, ``The somatic genomic landscape of glioblastoma,'' \textit{Cell}, vol. 155, no. 2, pp. 462--477, 2013.
\bibitem{mobadersany2018predicting} P. Mobadersany \textit{et al.}, ``Predicting cancer outcomes from histology and genomics using convolutional networks,'' \textit{Proc. Natl. Acad. Sci. U.S.A.}, vol. 115, no. 13, pp. E2970--E2979, 2018.
\bibitem{huang2020multimodal} Z. Huang \textit{et al.}, ``Multimodal deep learning for cancer prognosis,'' \textit{Nat. Commun.}, vol. 11, no. 1, p. 5740, 2020.
\bibitem{price2019privacy} B. A. Price and R. Cohen, ``Privacy and security in healthcare data management,'' \textit{Lancet Digit. Health}, vol. 1, no. 5, pp. e204--e212, 2019.
\bibitem{sheller2020federated} M. J. Sheller \textit{et al.}, ``Federated learning in medicine,'' \textit{Sci. Rep.}, vol. 10, no. 1, p. 12598, 2020.
\bibitem{mcmahan2017communication} B. McMahan \textit{et al.}, ``Communication-efficient learning of deep networks from decentralized data,'' in \textit{Proc. AISTATS}, 2017, pp. 1273--1282.
\bibitem{zhu2019deep} L. Zhu \textit{et al.}, ``Deep leakage from gradients,'' in \textit{Proc. NeurIPS}, 2019, pp. 14774--14784.
\bibitem{dwork2014algorithmic} C. Dwork and A. Roth, ``The algorithmic foundations of differential privacy,'' \textit{Found. Trends Theor. Comput. Sci.}, vol. 9, no. 3--4, pp. 211--407, 2014.
\bibitem{cox1972regression} D. R. Cox, ``Regression models and life-tables,'' \textit{J. R. Stat. Soc. Ser. B}, vol. 34, no. 2, pp. 187--202, 1972.
\bibitem{ishwaran2008random} H. Ishwaran \textit{et al.}, ``Random survival forests,'' \textit{Ann. Appl. Stat.}, vol. 2, no. 3, pp. 841--860, 2008.
\bibitem{katzman2018deepsurv} J. L. Katzman \textit{et al.}, ``DeepSurv: personalized treatment recommender system,'' \textit{BMC Med. Res. Methodol.}, vol. 18, no. 1, p. 24, 2018.
\bibitem{sheller2019multi} M. J. Sheller \textit{et al.}, ``Multi-institutional deep learning modeling without sharing patient data,'' in \textit{Proc. IEEE ISBI}, 2019, pp. 1173--1177.
\bibitem{pati2022federated} S. Pati \textit{et al.}, ``Federated learning enables big data for rare cancer boundary detection,'' \textit{Nat. Commun.}, vol. 13, no. 1, p. 7346, 2022.
\bibitem{stripelis2024federated} D. Stripelis \textit{et al.}, ``Federated learning on cross-institutional glioblastoma electronic health records,'' \textit{BMC Med. Inform. Decis. Mak.}, vol. 24, no. 1, p. 114, 2024.
\bibitem{abadi2016deep} M. Abadi \textit{et al.}, ``Deep learning with differential privacy,'' in \textit{Proc. ACM CCS}, 2016, pp. 308--318.
\bibitem{mironov2017renyi} I. Mironov, ``R\'enyi differential privacy,'' in \textit{Proc. IEEE CSF}, 2017, pp. 263--275.
\end{thebibliography}

\end{document}
